% Replacement proof for C.1.3; the parent owns sections/C.tex.
% Status may be changed to Proved.  No base-field exact sequence is used.
\begin{proof}
Put $V=H_{\B}(M)$, $f=f_{\infty,M}$ and
$G=\Aut^\otimes(H_{\B}|_{\langle M\rangle})\subset\GL(V)$.
Let $h:\mathbb S_{\R}\to\GL(V_{\R})$ be the Hodge homomorphism,
and let $H\subset\GL(V)$ be the smallest algebraic $\Q$-subgroup
whose scalar extension contains its image. This is the
Mumford--Tate group of $V$. It is connected: the connected group
$h(\mathbb S_{\R})$ contains the identity and is contained in
$H^0_{\R}$, while $H^0$ is defined over $\Q$; minimality gives
$H=H^0$.

The Hodge-pair identity gives
$fh(z)f^{-1}=h(\bar z)$. Since conjugation permutes
$h(\mathbb S_{\R})$, the rational element $f$ normalizes $H$.
As $f^2=1$, the finite union
\[
                         K=H\cup fH
\]
is a closed algebraic $\Q$-subgroup, with $K^0=H$.
The maps $h$ and $f$ act naturally on every tensor construction
and commute with every Hodge-pair morphism. Full faithfulness
in C.1.1 therefore gives $K\subset G$. In every mixed tensor
construction $T(V)$, that same full faithfulness gives
\[
 T(V)^G=\operatorname{Hom}(\mathbf1,T(V))
       =\{t\in T(V):t\text{ has type }(0,0),\ ft=t\}
       =T(V)^K.                                      \tag{C.1.3a}
\]
Here $H$-invariance is precisely invariance under $h$, by its
definition as the smallest algebraic $\Q$-group containing $h$;
for a rational vector, invariance under $h$ is the stated
Hodge-type condition. The assertion includes duals and direct
sums, so it determines Hom spaces between all tensor words.

For completeness, these invariants determine the group in this
setting without a normality assumption. Polarizable rational
Hodge structures are semisimple. If $W\subset X$ is a subobject
of Hodge pairs and $p:X\to W$ is a Hodge retraction, then
$(p+f_Wpf_X)/2$ is a Hodge retraction commuting with $f$.
It is a retraction because both summands restrict to $1_W$;
it commutes with $f$ by $f_X^2=f_W^2=1$. Thus the category of
polarizable Hodge pairs is semisimple, as is the tensor
subcategory in C.1.1. The corresponding $G$ is reductive.
The same semisimplicity of polarizable Hodge structures shows
that $H$ is reductive; its finite extension $K$ is linearly
reductive in characteristic zero, by averaging over $K/H$.
Every finite-dimensional representation of the closed subgroup
$K\subset\GL(V)$ is a subquotient of finite sums of mixed
tensor words in $V$, hence, by this reductivity, a direct
summand. Equation (C.1.3a) says that every $K$-equivariant map
between these words, including every idempotent, is
$G$-equivariant. Restriction $\operatorname{Rep}(G)\to
\operatorname{Rep}(K)$ is therefore fully faithful and
essentially surjective. Tensor reconstruction
\cite[Theorem 2.11]{DM1982} identifies the inclusion $K\hookrightarrow G$
with an isomorphism.

Consequently $G^0=H$ and $G=G^0\cup fG^0$. Its component
group is trivial or the constant group of order two. In the
second case left multiplication by $f\in G(\Q)$ is the
claimed isomorphism of $\Q$-schemes $G^0\to fG^0$.
For the homological category of C.1.2 its stated fullness
supplies the same equality (C.1.3a), and the same proof applies.
\end{proof}
